\section{Strengths, Limitations and Extensions}\label{sec:discussion}

\paragraph{Strengths.}
\begin{itemize}
  \item \textbf{Continuous-time physics, one mechanism per level.} L0--L3 are ODEs driven by a piecewise-constant
  usage input; switching one mechanism off (\cref{sec:structural}) measures what it is worth for TTE.
  \item \textbf{Closed forms} (F1--F8) explain and check the simulations: TTE at constant power within
  0.0009\,\% of isothermal L2, stranded charge within 0.17 points wherever a brown-out ends the run, approximate
  mean and variance under random usage.
  \item \textbf{Traceable and falsifiable.} Every parameter has a cited source, an open-data fit or a labeled
  assumption; one discharge constant is calibrated to phone data and the aging rates to the manufacturers'
  1000-cycle rating. The validation was pre-registered, and its failures are reported (31 of 59 criteria pass).
  \item \textbf{Three sources of uncertainty} (parameters: a 2048-member sample; model structure: variants of the
  main assumptions; usage randomness), plus Sobol' indices of 24 factor groups.
\end{itemize}

\paragraph{Limitations.}
\begin{itemize}
  \item \textbf{Cold capacity.} No equivalent circuit captures the transport-limited collapse at high C-rates
  (DFN: $+45$ to $+364\,\%$). At $-20\,^\circ$C L2 keeps 0.98 and 0.84 of the HG2 capacity (0.5C, 2C) where 0.73
  and 0.49 remain (activation energy 20.5 instead of 26.4\,kJ/mol). The brown-out size hinges on a pack resistance
  known within a factor of two (V7: 0.25 or 0.13\,$\Omega$). Cold TTEs are indicative.
  \item \textbf{One lumped thermal node} (A3): hot spots are quasi-static, and gaming time rests on two lumped
  parameters (enclosure factor and skin resistance, $S_{Ti}=0.40$ and 0.38).
  \item \textbf{Device floors and GNSS.} Android idle and video runtimes imply 1.5--2.8 times our power (TTE
  errors $+53$ to $+183\,\%$); a stationary iPhone's GPS draws about 0.03\,W, not our 0.27\,W prior.
  \item \textbf{Transferred laws.} $R\cdot Q$ scaling fails for a power cell (step resistance 21, not
  84\,m$\Omega$). The aging laws come from NMC and LFP cells \cite{SCH14,SCM18}, calibrated only to SOH 0.80 after
  1000 cycles; the fast-charge penalty rests on the LFP end of the plating prior, and the NASA test is only a
  consistency check of the capacity channel.
  \item \textbf{Assumed usage, optimistic defaults.} The CTMC rates are assumed, not fitted to logs, and carry
  44\,\% of the typical day's TTE variance. At 25\,$^\circ$C the median TTEs are 3--38\,\% shorter than the
  nominal ones (gaming's is 16\,\% longer).
  \item \textbf{No measured 4.45\,V OCV.} We stretch a 4.2\,V curve; other stretches move TTE by $\pm0.5\,\%$ at
  25\,$^\circ$C but by $-4$ to $+5\,\%$ on the cold, aged hike.
\end{itemize}

\paragraph{Extensions.}
(i)~An extended Kalman filter on $(\soc,v_1)$ with $R_0$ as a slow state would feed F4 with the measured
resistance and close the loop for the burst governor and gauge (\cref{sec:recs}). (ii)~A diffusion state in the
cold: the kinetic-battery variant X9 \cite{MM93} closes most of the 2C HG2 gap (0.52 vs.\ 0.49) and passes the
DFN test at 0\,$^\circ$C, 0.5C. (iii)~Two thermal nodes (processor hot spot, cell) would time throttling and cell
heating separately. (iv)~Learning the usage chain on the device: session rate and gaming probability carry most of
the usage-model variance (first-order indices 0.28, 0.45). (v)~Charging as optimal control over limit, rate and
timing, with plating as a constraint.
